\subsection{Smooth algebras}

Lemma 5.--- Let \(\rho: A \to B\) be a local homomorphism of Noetherian local rings. Suppose that B is essentially of finite type over A. For the A-algebra B to be formally smooth, it is necessary and sufficient that the A-module B be flat and that the \(\kappa_{A}\)-algebra \(\kappa_{A} \otimes_{A} B\) be absolutely regular.

There exists an integer \(n \geqslant 0\), a prime ideal q of \(A[T_{1}, \ldots, T_{n}]\), and a surjective homomorphism h from \(A[T_{1}, \ldots, T_{n}]_{q}\) to B. Let C denote the local A-algebra \(A[T_{1}, \ldots, T_{n}]_{q}\); it is formally smooth ( \(n^{\circ} 3\), Example 2 and \(n^{\circ} 2\), Prop. 4, a)) and flat over A, and one may identify B with the A-algebra C/J, where \(J = \text{Ker}(h)\).

Set \(\overline{\mathrm{C}} = \kappa_{\mathrm{A}} \otimes_{\mathrm{A}} \mathrm{C}\) and \(\overline{\mathrm{B}} = \kappa_{\mathrm{A}} \otimes_{\mathrm{A}} \mathrm{B}\). Suppose B is formally smooth over A. The \(\kappa_{\mathrm{A}}\)-algebra \(\overline{\mathrm{B}}\) is then formally smooth (no. 2, prop. 4, b)), hence absolutely regular (no. 5, cor. 2 of th. 2). Moreover, since \(\overline{\mathrm{C}} / \overline{\mathrm{JC}}\) is identified with \(\overline{\mathrm{B}}\) and since the \(\kappa_{\mathrm{A}}\)-algebra \(\overline{\mathrm{C}}\) is formally smooth, the ideal \(\overline{\mathrm{JC}}\) of \(\overline{\mathrm{C}}\) is completely secant (no. 5, th. 2). It then follows from § 5, no. 6, prop. 6 that the A-module B is flat.

Suppose conversely that B is flat over A and that the \(\kappa_{\mathrm{A}}\)-algebra \(\overline{\mathbf{B}}\) is absolutely regular. Then the \(\kappa_{\mathrm{A}}\)-local algebra \(\overline{\mathbf{B}}\) is formally smooth (Remark 1 of no. 9 with \(k = k_0 = \kappa_{\mathrm{A}}\)). Set \(\overline{\mathbf{J}} = \kappa_{\mathrm{A}} \otimes_{\mathrm{A}} \mathrm{J}\); since B is a flat A-module, the canonical map \(\overline{\mathbf{J}} \to \overline{\mathbf{J}}\overline{\mathbf{C}}\) is bijective and \(\overline{\mathbf{B}}\) is identified with \(\overline{\mathbf{C}}/\overline{\mathbf{J}}\). It follows (Remark 1 of no. 2) that the canonical map

\[
\overline {{\mathrm{J}}} / \overline {{\mathrm{J}}} ^ {2} \longrightarrow \overline {{\mathrm{B}}} \otimes_ {\overline {{\mathrm{C}}}} \Omega_ {\kappa_ {\mathrm{A}}} (\overline {{\mathrm{C}}})
\]

is injective and admits a retraction. Now \(\overline{\mathrm{J}}/\overline{\mathrm{J}}^{2}\) is identified with \(\kappa_{\mathrm{A}}\otimes_{\mathrm{A}}\mathrm{J}/\mathrm{J}^{2}\), hence with \(\overline{\mathrm{B}}\otimes_{\mathrm{B}}\mathrm{J}/\mathrm{J}^{2}\). On the other hand, the \(\overline{\mathrm{C}}\)-module \(\Omega_{\kappa_{\mathrm{A}}}(\overline{\mathrm{C}})\) is canonically isomorphic to \(\overline{\mathrm{C}}\otimes_{\mathrm{C}}\Omega_{\mathrm{A}}(\mathrm{C})\) (A, III, p. 136, prop. 20), hence \(\overline{\mathrm{B}}\otimes_{\overline{\mathrm{C}}}\Omega_{\kappa_{\mathrm{A}}}(\overline{\mathrm{C}})\) is canonically isomorphic to \(\overline{\mathrm{B}}\otimes_{\mathrm{C}}\Omega_{\mathrm{A}}(\mathrm{C})\). Passing to the quotient by the maximal ideal of \(\overline{\mathrm{B}}\), one obtains an injective homomorphism

\[
\kappa_ {B} \otimes_ {B} J / J ^ {2} \longrightarrow \kappa_ {B} \otimes_ {C} \Omega_ {A} (C)
\]

which is precisely \(\bar{d}_{\kappa_{\mathrm{B}}}\). Thus B is formally smooth over A (th. 3).

THEOREM 4.--- Let A be a Noetherian ring and B an A-algebra essentially of finite type. The following conditions are equivalent:

\begin{enumerate}
\def\labelenumi{(\roman{enumi})}
\item
  the A-algebra B is formally smooth;
\item
  for every \(\mathfrak{q} \in \operatorname{Spec}(\mathrm{B})\) , the A-algebra \(\mathrm{B}_{\mathfrak{q}}\) is formally smooth (resp. formally smooth for the topology \(\mathfrak{q}\mathrm{B}_{\mathfrak{q}}\) -adic);
\item
  the A-module B is flat and, for every \(\mathfrak{p} \in \operatorname{Spec}(A)\) , the \(\kappa(\mathfrak{p})\) -algebra \(\kappa(\mathfrak{p}) \otimes_{A} B\) is absolutely regular;
\item
  the A-module B is flat and, for every regular A-algebra R, the ring R ⊗\_A B is regular;
\item
  the A-module B is flat and the kernel of the homomorphism \(\mu:B\otimes_{A}B\to B\) defined by \(\mu(x\otimes y)=xy\) is a completely secant ideal.
\end{enumerate}

The equivalence of (i) and (ii) follows from cor. 2 of th. 3 (no. 9).

(i)⇒(v): suppose the A-algebra B is formally smooth. Let q be a prime ideal of B, and let p be its inverse image in A. The \(A_{\mathfrak{p}}\)-algebra \(B_{\mathfrak{q}}\) is formally smooth (prop. 4, a) of no. 2), hence flat (Lemma 5); consequently the A-module B is flat (II, § 3, no. 4, prop. 15). On the other hand, the ring B ⊗A B is Noetherian (§ 6, no. 1, cor. of prop. 2), so the ideal Ker μ is completely secant by cor. 3 of th. 2 (no. 5).

(v)⇒(iii): suppose condition (v) is satisfied. Set I = Ker(μ). Let p ∈ Spec(A). The map

\[
1 \otimes \mu : \kappa (\mathfrak {p}) \otimes_ {\mathrm{A}} (\mathrm{B} \otimes_ {\mathrm{A}} \mathrm{B}) \rightarrow \kappa (\mathfrak {p}) \otimes_ {\mathrm{A}} \mathrm{B}
\]

is identified with the map

\[
\mu_ {\mathfrak {p}}: (\kappa (\mathfrak {p}) \otimes_ {\mathrm{A}} \mathrm{B}) \otimes_ {\kappa (\mathfrak {p})} (\kappa (\mathfrak {p}) \otimes_ {\mathrm{A}} \mathrm{B}) \rightarrow \kappa (\mathfrak {p}) \otimes_ {\mathrm{A}} \mathrm{B}
\]

deduced from the multiplication of the \(\kappa(\mathfrak{p})\) -algebra \(\kappa(\mathfrak{p})\otimes_{\mathrm{A}}\mathrm{B}\) . The ideal \(\operatorname{Ker}(\mu_{\mathfrak{p}})\) is identified with \(\mathrm{I}(\kappa(\mathfrak{p})\otimes_{\mathrm{A}}(\mathrm{B}\otimes_{\mathrm{A}}\mathrm{B}))\) . It is completely secant since the A-module B is flat ( \(\S 5\) , no. 6, prop. 6). Assertion (iii) then follows from prop. 8 of \(\S 6\) , no. 5.

(iii)⇒(ii): let q be a prime ideal of B, and p its inverse image in A. Under the hypotheses of (iii), the A \(_{p}\) -module B \(_{q}\) is flat, and the κ(p)-algebra κ(p) ⊗A \(_{p}\) B \(_{q}\) , which is identified with a ring of fractions of κ(p) ⊗A B, is absolutely regular (§ 6, no. 4, prop. 6). It follows from lemma 5 that B \(_{q}\) is formally smooth over A \(_{p}\) , hence over A (no. 2, prop. 3 and 4).

(iii)⇒(iv): let us place ourselves under the hypotheses of (iii). Let R be a regular A-algebra. The R-module R ⊗A B is flat (I, § 2, no. 7, cor. 2 to prop. 8). Let r be a prime ideal of R and p its inverse image in A; the ring κ(r)⊗R(R⊗AB), which is identified with κ(r)⊗κ(p) (κ(p)⊗A B), is regular (§ 6, no. 4, cor. 2 of prop. 7). The ring R ⊗A B is therefore regular (§ 4, no. 5, cor. of prop. 9).

(iv)⇒ (iii): let p be a prime ideal of A and let k be an extension of κ(p); under the hypotheses of (iv), the ring \(k \otimes_{\kappa(\mathfrak{p})} (\kappa(\mathfrak{p}) \otimes_{\mathrm{A}} \mathrm{B})\) , which is identified with \(k \otimes_{A} B\) , is regular, whence (iii).

DEFINITION 2.--- Let A be a Noetherian ring. An A-algebra B is said to be smooth if it is essentially of finite type and if it satisfies the equivalent conditions of Theorem 4.

PROPOSITION 10.--- Let A be a Noetherian ring.

\begin{enumerate}
\def\labelenumi{\alph{enumi})}
\item
  Let A' be a Noetherian A-algebra and B a smooth A-algebra. Then the A'-algebra A' ⊗\_A B is smooth.
\item
  Let B be a smooth A-algebra and C a smooth B-algebra. Then the A-algebra C is smooth.
\item
  Let B and C be two smooth A-algebras. Then the A-algebra B ⊗A C is smooth.
\end{enumerate}

This follows from prop. 4 of no. 2 and the analogous statements for algebras essentially of finite type (§ 6, no. 1).

Examples.--- 1) Smooth algebras over a field \(k\) are precisely the essentially finite-type and absolutely regular \(k\)-algebras.

\begin{enumerate}
\def\labelenumi{\arabic{enumi})}
\setcounter{enumi}{1}
\tightlist
\item
  Let A be a Noetherian ring, \(\mathbf{T} = (\mathrm{T}_i)_{i\in I}\) a finite family of indeterminates. The A-algebra A{[}T{]} is smooth. More generally, let \(\mathrm{F}_1,\ldots ,\mathrm{F}_m\) of elements of A{[}T{]}, and let B be the A-algebra A{[}T{]}/ \((\mathrm{F}_1,\dots ,\mathrm{F}_m)\). If in every maximal ideal \(\mathfrak{n}\) of B the class (mod. \(\mathfrak{n}\) ) of the matrix \(\left(\frac{\partial\mathrm{F}_j}{\partial\mathrm{T}_i}\right)\) is of rank \(m\) , the A-algebra B is smooth (th. 3 of no. 9).
\end{enumerate}

